\begin{abstract}
  While parameter-efficient fine-tuning (PEFT) enabling modular adapters (e.g., LoRA) allows cost-effective multi-task serving, dynamic model-merging routers often rely on empirical heuristics that lack learning-theoretic guarantees and collapse under heteroscedastic domain noise. In this work, we present \textbf{PAC-ZCA}, a mathematically rigorous, learning-theoretic framework for dynamic model merging that optimizes a strictly temperature-only Gibbs routing policy using parameter-space PAC-Bayesian generalization bounds. We model sample-wise routing as a randomized ensembling policy over unsupervised, task-agnostic \textbf{Subspace Energy Projection (SEP)} features. Under a Gaussian prior/posterior over the log-temperature parameters, we establish a clean mathematical duality between minimizing the parameter-space KL-divergence (which limits the router's Lipschitz constant) and maximizing output routing entropy (which encourages smooth activation blending). Rather than relying on heuristic sweeps, PAC-ZCA solves for the optimal, task-specific temperatures by directly minimizing our derived PAC-Bayesian bound on a tiny, offline calibration set (16 samples per task). Evaluated inside a challenging 14-layer Coordinate Sandbox with extreme heteroscedastic noise, we employ disjoint calibration splits to resolve the double data-dependency flaw under SVD, fully satisfying McAllester's theorem. Under this sound split, PAC-ZCA (Block) achieves \textbf{64.16\% $\pm$ 2.23\%} joint classification accuracy, outperforming raw-coordinate SABLE by \textbf{+23.70\%} and matching the mean performance of standard unregularized Empirical Risk Minimization while successfully reducing ensembling variance. Furthermore, on real images (MNIST, Fashion-MNIST, CIFAR-10) using a pre-trained ResNet-18 backbone, PAC-ZCA achieves \textbf{70.87\% $\pm$ 2.20\%} joint accuracy, outperforming standard SABLE (65.67\%) and strictly outperforming standard unregularized ERM (69.47\%) while maintaining highly stable ensembling performance, demonstrating the practical power of learning-theoretic ensembling on real-world edge serves.
\end{abstract}
